\section{本章のまとめ}
\label{sec:quantum-basics-summary}

本章では，量子機械学習を理解するために必要な量子計算の基礎を説明した．
量子ビットの純粋状態は，計算基底$\ket{0}$と$\ket{1}$の複素線形結合として表され，
測定結果の確率は確率振幅の絶対値の二乗によって与えられる．
複数量子ビットの状態空間はテンソル積で構成され，積状態として分解できない状態として
量子もつれが存在する．

量子ゲートはユニタリ変換であり，量子回路は複数の量子ゲートを組み合わせて量子状態を変化させる．
量子測定では確率的に結果が得られ，観測量の期待値は量子状態とエルミート演算子を用いて表される．
また，古典データを量子回路へ入力する方法として，基底符号化と振幅符号化を説明した．

以上により，量子状態の表現，量子ゲートによる状態変化，測定による情報の取り出し，
および古典データを量子状態へ対応させる方法という，量子計算の基本事項を整理した．
